% Models and training: the same sequence of models for every task, so each
% model paper inputs this file under its own heading.
% The prose follows docs/writing.md.
Every task is given the same sequence of models, so that each step of the
sequence measures the contribution of one component. A model is fitted on the
training set, selected on the validation set and scored once on the test set.
Table~\ref{tab:runs} lists the models of this task.

\paragraph{Constant baseline.}
\run{constant} predicts the same value for every example: the mean of the
training targets (for a binary target, the base rate), or their median when
the primary metric is an absolute error. A useful model has to improve on it.

\paragraph{Gradient-boosted decision trees.}
\ifcontextgbdt
\run{gbm-context} is a gradient-boosted decision tree model (GBDT), trained
with LightGBM~\cite{lightgbm} on the scalar context features alone. It never
sees the board, and so measures how much of the target the context already
determines. \run{gbm-hand} adds the hand-crafted features.
\else
\run{gbm-hand} is a gradient-boosted decision tree model (GBDT), trained with
LightGBM~\cite{lightgbm} on the scalar context features and the hand-crafted
features.
\fi
We refer to \run{gbm-hand} as \emph{the GBDT baseline}: every other model is
compared with it on the same test examples. GBDT is trained for at most
\V{method.gbm.rounds} boosting rounds with up to \V{method.gbm.leaves} leaves
per tree and a learning rate of \V{method.gbm.rate}. Each tree is fitted on a
random \V{method.gbm.rows.percent}\% of the examples and
\V{method.gbm.features.percent}\% of the features, and boosting stops when the
validation score has not improved for \V{method.gbm.patience} rounds.

\paragraph{Neural networks.}
A network consists of an encoder, which maps the square tokens, the auxiliary
planes and the context features to a vector, and a prediction head, a
two-layer perceptron that maps this vector to the output. We use two encoder
architectures, so that the results show how much the architecture matters for
the task.

The \emph{transformer} encoder~\cite{attention} treats each square as a token.
The embedding of a token is the sum of a learned embedding of its piece, a
learned embedding of its square and a linear projection of the square's
auxiliary planes. The context features are projected by a small perceptron
and enter as one additional token. Every token attends to every other token
in every layer, so distant squares interact from the first layer on. To give
attention the geometry of the board, each attention head adds a learned bias
that depends on the file and rank offset between two
squares~\cite{relative}: files, ranks, diagonals and knight moves all
correspond to fixed offsets. The layers use pre-layer normalisation, and the
encoder outputs the context token concatenated with the mean of the 64 square
tokens.

The \emph{residual convolutional network} (ResNet) is a stack of residual
blocks of $3\times3$ convolutions over the $8\times8$ board, the architecture
that AlphaZero introduced for chess~\cite{alphazero}. Its input is one binary
plane per token value together with the auxiliary planes. A linear projection
of the context features is added at every square after the first convolution,
and the encoder outputs the mean and the maximum over the squares of the
final feature map. Unlike attention, a convolution relates only neighbouring
squares, so the receptive field grows with depth.

\ifsequence
\paragraph{Sequence model.}
An example of this task is a sequence of positions, so a network has two
stages below its prediction head. One of the two encoders above is applied to
every position of the sequence, with the position's per-move features in
place of the context features, and produces one vector per position. A
two-layer transformer then operates on these vectors in the order in which
the moves were played: each vector is a token, a learned positional embedding
is added to it, and the game-level context features enter as one additional
token. The prediction head receives that token concatenated with the mean of
the position tokens. Sequences in a batch are padded to the length of the
longest, and padded elements are masked: they are not encoded and no token
attends to them, so a prediction does not depend on the other examples in
the batch.

Model sizes are hyperparameters. \run{conv-small} is small enough to train
on a CPU. An example consists of tens of positions, so
\run{transformer-small} and \run{transformer-medium} are trained on a GPU.
\run{transformer-medium-hand} is \run{transformer-medium} with the
hand-crafted features added to its context features.
Table~\ref{tab:settings} in the appendix lists the hyperparameters.
\else
Model sizes are hyperparameters. \run{conv-small} and \run{transformer-small}
are small enough to train on a CPU. \run{transformer-medium} is trained on a
GPU, and \run{transformer-medium-hand} is the same network with the
hand-crafted features added to its context features.
Table~\ref{tab:settings} in the appendix lists the hyperparameters.
\fi

\paragraph{Training.}
All networks are trained in the same way: with AdamW~\cite{adamw}, a learning
rate that increases linearly over the first \V{networks.warmup.percent}\% of
the optimisation steps and then follows a cosine schedule to zero, gradient
clipping, and mixed precision on a GPU. The loss matches the primary metric:
mean squared error for a target scored by RMSE, a Huber loss for a target
scored by MAE, and binary cross-entropy for a probability. A regression
target is standardised with the mean and standard deviation of the training
set. After every epoch the network is evaluated on the validation set. We
keep the weights of the best epoch, and stop training when
\V{networks.patience} consecutive epochs have not improved on it.
\ifmore
A model trained on a larger training set is given proportionally fewer
epochs, and a correspondingly shorter patience
(Section~\ref{sec:ablations}).
\fi

\begin{table}[htbp]\centering\footnotesize
\block{runs}
\caption{The models, best first. \emph{Selected} is the number of trees or
the training epoch chosen on the validation set. \emph{Time} is the total
wall-clock time for training and evaluation.}
\label{tab:runs}
\end{table}
